\section{Purpose}

\textbf{[STIPULATION]} $\xi = 4/30000$ is a stipulation (A020). This document
answers whether this number is \emph{only} empirically anchored, or whether it
follows from the structure of the system. The result: $\xi$ is not arbitrary.
The exponent $-4$ in the order of magnitude $10^{-4}$ follows from the fractal
spacetime dimensionality $D_f=3-\xi$ and the four-dimensionality of the universe.
The mass-scaling exponent $\kappa=7$ is a \emph{choice of the decomposition}
$R_{pe} = K\cdot(4/3)^7$ with remainder term $K$; there is no Casimir
confirmation of the factor $4/3$.
The order of magnitude of $\xi$ is thus anchored -- but \emph{not} derivable from the three stipulations of A020
alone; a measured mass anchor remains necessary.

\section{The Circularity Problem}

\textbf{[S]} The original presentation appears circular:
\[
  \frac{m_p}{m_e} = 245\times\Bigl(\frac43\Bigr)^7
  \;\Rightarrow\; \xi = \frac{4}{30000}.
\]
\textbf{Objection:} Why $\kappa=7$? Why $K=245$? Is this not backward fitting?

\textbf{[B]} The answer: the exponent $\kappa=7$ is a choice of the decomposition
$R_{pe} = K\cdot(4/3)^\kappa$ with remainder term $K$ -- with free $K$, every $\kappa$ fits exactly.
That is the core of the document; the rest demonstrates it.

\section{The e-p-$\mu$ System as Proof}

\textbf{[K]} The three fundamental mass ratios:
\begin{align}
  R_{pe}    &= \frac{m_p}{m_e}   = 1836{.}15267343, \notag\\
  R_{\mu e} &= \frac{m_\mu}{m_e} = 206{.}7682830, \notag\\
  R_{p\mu}  &= \frac{m_p}{m_\mu} = 8{.}880. \notag
\end{align}
Multiplicativity forces $R_{pe} = R_{\mu e}\times R_{p\mu}$. For the ansatz
$R_{pe} = K\times(4/3)^\kappa$ one need only test the exponent:

\begin{center}
\renewcommand{\arraystretch}{1.3}
{\small
\begin{tabular}{rccr}
\toprule
$\kappa$ & Prediction $R_{pe}$ & Consistent & Error \\
\midrule
6 & $245\times(4/3)^6 = 1376{.}6$ & \texttimes & 25{.}0\,\% \\
\textbf{7} & $\mathbf{245\times(4/3)^7 = 1835{.}4}$ & \checkmark & \textbf{0{.}04\,\%} \\
8 & $245\times(4/3)^8 = 2447{.}2$ & \texttimes & 33{.}3\,\% \\
\bottomrule
\end{tabular}
}
\renewcommand{\arraystretch}{1.0}
\end{center}

\textbf{[K]} The table keeps $K=245$ (fitted to $\kappa=7$) fixed and can therefore
single out only $\kappa=7$. For every $\kappa$, $K_\kappa = R_{pe}/(4/3)^\kappa$
fits exactly: $\kappa = 6,\ 7,\ 8$ give $K = 326.8;\ 245.1;\ 183.8$. The formula
\[
  \kappa = \frac{\ln(R_{pe}/K)}{\ln(4/3)}
  = \frac{\ln(1836{.}15/245)}{\ln(1{.}3333)} \approx 7{.}000
\]
presupposes $K = 245$. $\kappa=7$ is a choice of the decomposition, not a
selection; the table shows no uniqueness.

\section{The Fundamental Justification of the $10^{-4}$ Scale}

\textbf{[B]} Why $10^{-4}$? The argument has three steps:

\textbf{(1) From the fractal dimension.} The fractal dimension
$D_f=3-\xi\approx2{.}9998667$ generates a discrete scale hierarchy. The
relative deviation from the integer dimension:
\[
  \delta = 1-\frac{D_f}{3} = 1-\frac{3-\xi}{3} = \frac{\xi}{3}
  \;\approx\; 4{.}44\times10^{-5},
\]
so $\xi=3\delta=3-D_f$; $\xi$ is the absolute deviation of the fractal
dimension from~3.

\textbf{(2) From the spacetime dimensionality.} In $d$-dimensional spaces
one expects natural scalings $\xi_d\sim(10^{-1})^d$. For $d=4$:
\[
  \xi_4\sim(10^{-1})^4 = 10^{-4}.
\]
This explains the order of magnitude; the prefactor $4/3$ does not follow
from it (for step (3) see there).

\textbf{(3) From the Casimir effect.} The Casimir energy per area between
ideal parallel plates is
\[
  \frac{E_\text{Casimir}}{A} = -\frac{\pi^2\hbar c}{720\,a^3}
\]
without a factor $4/3$; there is no Casimir confirmation of the factor $4/3$.
Only the fourth ($4/3$) -- not the fifth ($3/2$) or third ($5/4$) -- gives
consistent results for the mass spectrum. Consistency over five independent
ratios:

\begin{center}
\renewcommand{\arraystretch}{1.3}
{\small
\begin{tabular}{lccc}
\toprule
Ratio & Experiment & T0 with $\kappa=7$ & Error \\
\midrule
$m_p/m_e$       & 1836{.}1527 & 1835{.}4 & 0{.}04\,\% \\
$m_\mu/m_e$     & 206{.}7683  & input & -- \\
$m_p/m_\mu$     & 8{.}880     & 8{.}877 & 0{.}04\,\% \\
$m_\tau/m_\mu$  & 16{.}817    & input & -- \\
$m_n/m_p$       & 1{.}001378  & 1{.}001333 & 0{.}004\,\% \\
\bottomrule
\end{tabular}
}
\renewcommand{\arraystretch}{1.0}
\end{center}

\section{What $K = 245$ Is and What It Is Used For}

\textbf{[K]} $K=245$ is a \emph{remainder term}, not a fundamental object.
$\ln(R_{pe})/\ln(4/3) = 26{.}12$ -- not an integer. To obtain an integer
exponent $\kappa=7$, one decomposes:
\[
  K = \frac{R_{pe}}{(4/3)^7}
  = \frac{1836{.}152}{7{.}492} \approx 245{.}1.
\]
$K$ names the remainder after extracting $(4/3)^7$. It appears in no other
FFGFT formula. If $R_{pe}$ were exactly $245\times(4/3)^7=1835{.}43$,
$K$ would be exactly an integer; the small remainder $R_{pe}-245\times(4/3)^7=0{.}73$
produces the 0{.}04\,\% deviation.

\textbf{[S]} The prime factorisation $245=5\times7^2$ is a post-hoc
observation -- $7$ appears because of $\kappa=7$, $5$ is coincidental.
There is no geometric derivation of $K=245$; a corresponding formula
($245\approx\phi^{12}/(1-\xi)^2$) from the legacy corpus is numerically
wrong ($\phi^{12}/(1-\xi)^2\approx322$, deviation 31\,\%).

\textbf{[K]} \textbf{Honest assessment.} The anchors via fractal dimension (1)
and spacetime dimensionality (2) explain the order of magnitude $10^{-4}$.
There is no Casimir confirmation of the factor $4/3$ (step (3)). All arguments
hang on the e-p-$\mu$ system as an empirical anchor. $\kappa=7$ is a choice of
the decomposition, not a selection. There is no derivation of $\xi$ from the
three stipulations of A020 alone (without measured masses).

\section{Verification Script}

\begin{itemize}[leftmargin=2em,itemsep=2pt,topsep=2pt]
  \item Numerical verification: $\kappa=7$ as the unique consistent solution
    for $R_{pe}$, $R_{\mu e}$, $R_{p\mu}$; errors for $\kappa=6,7,8$:
    \texttt{python/A\_Serie\_Skripte/a015\_xi\_ursprung.py}
    (An error comparison at fixed $K=245$ does not show uniqueness of $\kappa$.)
\end{itemize}

\section{Open Questions}

\textbf{There will be no proof of $\xi$ -- and that is not a deficiency.}
$\xi$ is a fundamental parameter of the theory -- analogous to $\alpha$ in
the Standard Model or $G$ in gravitation. Fundamental parameters are not
proved from deeper principles; they are motivated by plausibility arguments,
supported by consistency, and confirmed by measurements. A ``purely geometric
derivation without empirical input'' is in principle impossible for a
fundamental parameter -- this holds for any theory.

\textbf{What exists: plausible reasons at several levels.}

\textbf{(1) The fourth as motivated basis.} Of all harmonic
intervals -- fourth $4/3$, fifth $3/2$, third $5/4$ -- only the fourth
yields a simple near-integer remainder $K \approx 245$ for $\kappa=7$.
With a free remainder term $K$, every base fits exactly; singling out $4/3$
is a choice, not an experimental selection mechanism.

\textbf{(2) Fractal deviation.} $\xi=3-D_f$ is the absolute deviation
of the fractal dimension $D_f=3-\xi$ from the integer dimension~3.
The equation $\xi=\xi$ is circular, but the consistency ($D_f$ and $\xi$
determine each other uniquely) is a structural argument.

\textbf{(3) Natural scaling.} In $d$-dimensional spaces $\xi_d\sim(10^{-1})^d$
is the natural scaling; for $d=4$ (3 spatial + 1 time) it follows that
$\xi_4\sim10^{-4}$.

\textbf{(4) Prime factor structure.} $\xi=4/30000=2^2/(3\cdot2^4\cdot5^4)
=1/(3\cdot2^2\cdot5^4)$ -- factor~3 corresponds to the spatial dimensions,
factor~4 to the spacetime dimensions, $5^4$ to the fractal structure.

\textbf{(5) Music theory analogue.} Pythagorean 3:4:5 tetrahedra as a
structural parallel to the harmonic ratios -- not a proof, but a methodological
origin (accordion building, acoustics, resonance physics) that suggests the
interval 4:3 as a natural unit.

Together these reasons motivate $\xi=4/30000$ as a plausible, consistent
value -- not as a proved one. This is the honest state, and it is the same
state as for every other fundamental parameter in physics.

\textbf{$K=245$ is a remainder term, not a fundamental parameter.} It follows
from $R_{pe}$ and $\kappa=7$ -- an intermediate value, not an independent
quantity (A015, section on K=245).

\section{Supersedes}

This document subsumes: Doc.~009 ($\xi$-origin: circularity problem, e-p-$\mu$
system, Casimir confirmation (not tenable, see step (3) of the $10^{-4}$ justification), $K=245$, complete consistency system), Doc.~180
(complete derivation chain in five steps: T0 Lagrangian $\to$ self-consistency
condition $\to$ uniqueness of $\xi$ $\to$ $G$ $\to$ $L_0$, connection to the
SI system).

\section{Use of AI Tools}

This document was created with the assistance of AI tools. The questions,
stipulations, and all substantive decisions come from the author; the tools
formulate, compute, and create verification scripts.
Responsibility for content and errors rests entirely with the author.
Full wording of the rules in A010.
